\documentclass[11pt,letterpaper]{article}
\usepackage[T1]{fontenc}
\usepackage[utf8]{inputenc}
\usepackage{lmodern,microtype}
\usepackage[margin=.91in,headheight=15pt]{geometry}
\usepackage{amsmath,amssymb,amsthm,mathtools}
\usepackage{booktabs,longtable,tabularx,array,enumitem}
\usepackage{xcolor,fancyhdr,listings,xurl,needspace}
\usepackage[colorlinks=true,linkcolor=blue!45!black,citecolor=blue!45!black,urlcolor=blue!45!black]{hyperref}
\definecolor{ink}{RGB}{25,48,65}
\definecolor{pale}{RGB}{246,248,250}
\setlength{\emergencystretch}{3em}
\setlength{\parskip}{.25em}
\setlist{itemsep=3pt,topsep=5pt,leftmargin=1.7em}
\newtheorem{theorem}{Theorem}[section]
\newtheorem{lemma}[theorem]{Lemma}
\newtheorem{proposition}[theorem]{Proposition}
\newtheorem{corollary}[theorem]{Corollary}
\theoremstyle{definition}\newtheorem{definition}[theorem]{Definition}
\theoremstyle{remark}\newtheorem{remark}[theorem]{Remark}
\newcommand{\ipra}{\mathsf{iPRA}}
\newcommand{\CG}{\mathcal G}
\newcommand{\UR}{\mathcal U_{\ne0}}
\newcommand{\NC}{\mathcal N}
\newcommand{\FV}{\operatorname{FV}}
\newcommand{\file}[1]{{\small\nolinkurl{#1}}}
\newcommand{\equider}{\mathrel{\dashv\vdash}}
\lstset{basicstyle=\ttfamily\footnotesize,columns=fullflexible,keepspaces=true,
 breaklines=true,showstringspaces=false,frame=single,rulecolor=\color{black!18},
 backgroundcolor=\color{pale},aboveskip=8pt,belowskip=8pt}
\pagestyle{fancy}\fancyhf{}
\fancyhead[L]{\small A guarded orbit presentation of iPRA}
\fancyhead[R]{\small November 11, 2026}
\fancyfoot[C]{\thepage}
\hypersetup{pdftitle={A Guarded Orbit Presentation of iPRA},
 pdfauthor={Research report prepared for Vladimir Reshetnikov}}
\title{\color{ink}\bfseries A Guarded Orbit Presentation of iPRA\\[8pt]
\large A further reduction of the ConciseIPRA checker,\\with an explicit equivalence boundary}
\author{Research report prepared for Vladimir Reshetnikov}
\date{November 11, 2026}
\begin{document}
\maketitle
\begin{abstract}
We give a further, small reduction of the current ConciseIPRA presentation.
Elimination selects its major premise from its source term, and introduction
selects cases or equation progress from the root of its matrix. These two
reductions already exist separately. They do not compose inside the strictly
guarded calculus: an assumption $0=t$ cannot there yield $t=0$.
The new argument repairs that failure using the intermediate orbit formula
$D_{r,y}(y=I(y,t;u.Su))$. It supplies the one closure needed to recover
zero-source equality symmetry. Stable-left numerical cases then suffice to
complete the bootstrap, and the original guarded formula language is retained
for public statements. We prove that this presentation is equiconsistent and
whole-domain bi-interpretable with full first-order intuitionistic primitive
recursive arithmetic, with actual equality and identity numerical composites.

Three complete Python programs improve the separate attained coordinates
from 259 tokens, 283 AST nodes and 526 normalized bytes to 258, 281 and 525.
These are different programs. The compact comparison assumes canonical
public inputs. A supplied checker that also validates the public grammar is
1205 bytes, compared with 1208 for the prior guarded interface and 1109 for
the prior scalar interface. Thus this is a small positive result under the
repository's canonical-input convention, not an absolute simplest theory or
a new minimum for every strict interface. The new metatheorems are written
proofs; no new Lean or Rocq verification is claimed.
\end{abstract}

\begin{center}
\fcolorbox{ink}{pale}{\begin{minipage}{.91\linewidth}
\textbf{The mathematical contribution.} Intermediate formulas can be more
expressive than the public language without changing its theory. Here a
single family of iterator-orbit formulas makes two otherwise incompatible
rule reductions work together. The statement of the theory includes this
distinction; it is not an omitted validation cost.
\end{minipage}}
\end{center}
\newpage
\tableofcontents
\newpage

\section{Result, reference point, and meaning of simpler}

The reference is the user's ProveIt repository at commit
\nolinkurl{1517d9d79ce33f631e0f82960c849427cbba23bf}, under
\file{Logic/PrimitiveRecursiveArithmetic/ConciseIPRA} \cite{repo}.
At that commit the Lean kernel already has equational transport and direct
$Q$ progress; literal-zero source elimination is also proved in Lean.
The written packages additionally contain guarded, root-determined
introduction. None of these individual changes is claimed as new here.
The current README lists twenty-eight research packages, and reports
separate checker minima for the \file{source_dispatch} family.

The new result is the \emph{compatibility proof with administrative orbits}.
It is necessary: simply intersecting the two reduced rule systems does not
preserve equality reasoning. The obstruction and its repair are proved
below. The numerical syntax, scalar arithmetic, logical encodings, and
existing iPRA bridge are inherited and explicitly attributed.

\begin{theorem}[Public theory and numerical composites]\label{thm:main}
Let $\CG$ be the public language of Definition~\ref{def:public}, and permit
proofs in the administrative calculus $\UR$ of
Section~\ref{sec:rules}. For every finite context $\Gamma$ and conclusion $A$
in $\CG$, this calculus derives exactly the sequents of the inherited
guarded presentation. There are effective interpretations between its
ordinary realization and full first-order $\ipra$, and each formula is
provably equivalent to its round trip. Both numerical domains are whole,
equality is actual equality, and both numerical composites have graph $x=y$.
In particular, the two theories are equiconsistent.
\end{theorem}

Here ``full first-order'' describes the available logic, including arbitrary
quantified formulas. It does not mean full induction. On the iPRA side,
induction is restricted to quantifier-free formulas; the arithmetic progress
rule below applies only to literal equations.

\paragraph{The comparison is deliberately qualified.}
There is no proved universal lower bound of 259 tokens, 283 nodes, or
526 bytes. Those are attained coordinates of inspected programs.
The smaller programs below recognize exactly the intended theorems on
canonical public inputs, while allowing broader intermediate syntax.
The cost of validating canonical inputs is reported separately and also
included in one standalone strict program. The latter does not beat the
1109-byte old scalar interface. A claim of an unqualified new global minimum
would therefore be unsupported.

\section{The complete presentation}\label{sec:rules}

\subsection{Numerical terms and finite descriptions}

There is one numerical sort. Terms are finite expressions
\[
 t ::= x\mid0\mid S(t)\mid I(n,a;u.h).
\]
Only $u$ is bound in $h$. The counter $n$ and initializer $a$ are outside
that binder. Other variables in $h$ are passive parameters. Substitution
is capture-free, and bound variables may be renamed.
Write $1=S0$, $2=S1$, and
\[
 C(v,a,b):=I(v,a;u.b),
\]
with $u$ fresh for $b$. These are abbreviations, not extra constructors.
The finite descriptions can equivalently be expanded to ordinary first-order
function symbols, one for each fixed previously constructed iterator body.
No variable ranges over functions or over descriptions, and no uniform
total evaluator is postulated internally.

\subsection{Two explicit formula grammars}

The administrative grammar is
\[
 A ::= (s=t)\mid Q_y A\mid D_{x,y}A.
\]
$Q$ binds $y$; $D$ binds the distinct variables $x,y$. A \emph{literal
equation} is a formula whose outermost constructor is equality.
Every other administrative formula is called compound. Retained contexts
are finite reusable lists of administrative formulas.

\begin{definition}[Public formulas]\label{def:public}
The public language $\CG$ is the least grammar satisfying:
\begin{enumerate}
\item Every equation $s=t$ belongs to $\CG$.
\item $D_{x,y}B$ belongs to $\CG$ when $B$ belongs to $\CG$ and is compound.
\item An atomic matrix immediately below $D$ is allowed exactly when it has
      the shape
      \[
       D_{x,y}\bigl(I(y,0;u.h)=I(y,0;u.k)\bigr).
      \]
      The two state binders are independent, and $h,k$ may contain their
      state, $x,y$, and passive parameters.
\end{enumerate}
There is no public $Q$ constructor. It remains available in proofs.
The designated contradiction is the public equation $0=1$.
\end{definition}

In particular, the public language is exactly the earlier guarded language,
not a newly enlarged language with an additional orbit clause.
The orbit formulas used later occur only inside proofs. Every public guarded
atomic matrix has a provable zero instance, since both sides compute to zero.
This syntactic fact is the reason a single introduction scheme suffices.

\subsection{Axioms and inference rules}

The two computation schemas are available under every context:
\[
 I(0,a;u.h)=a,\qquad
 I(Sn,a;u.h)=h[I(n,a;u.h)/u].                       \tag{I0,IS}
\]
The remaining rules are precisely the following. No unlisted logical,
equality, successor, or induction axiom is primitive.

\paragraph{Hypothesis.}
If $A\in\Gamma$, then $\Gamma\vdash A$.

\paragraph{Source-selected elimination.}
For an administrative formula template $F(v)$, define
\[
 M(s,t,F)=
 \begin{cases}
 Q_vF(v),&s\text{ is literally }0,\\
 s=t,&s\text{ is not literally }0.
 \end{cases}
\]
The rule is
\[
 \frac{\Gamma\vdash M(s,t,F)\qquad\Gamma\vdash F(s)}
      {\Gamma\vdash F(t)}.                         \tag{E}
\]
Thus literal zero selects $Q$ elimination and every other syntactic source
selects equality transport. An iterator term denoting zero is a permitted
equality source. This is a test of syntax, never a test of its numerical value.

\paragraph{$D$ elimination.}
\[
 \frac{\Gamma\vdash D_{x,y}B(x,y)\qquad
       \Gamma,Q_yB(b,y)\vdash A}
      {\Gamma\vdash A}.                            \tag{DE}
\]
The eigenvariable $b$ is fresh for $\Gamma,A$ and the passive parameters of
$B$. Its local assumption is discharged.

\paragraph{Root-determined $D$ introduction.}
Put
\[
 \eta_B(z)=
 \begin{cases}z,&B\text{ is a literal equation},\\
               0,&B\text{ is compound}.
 \end{cases}
\]
Then
\[
 \frac{\Gamma,B(t,\eta_B(z))\vdash B(t,Sz)}
      {\Gamma\vdash D_{x,y}B(x,y)}.                 \tag{DI}
\]
Here $z$ is fresh for $\Gamma,t$ and the passive parameters; the local
assumption is discharged. The witness $t$ is any numerical term and is
chosen before introducing $z$. The certificate has no separate
case/progress choice. At an equation the rule is progress; at a compound
matrix it is cases. Arbitrary unguarded atomic matrices are allowed as
intermediate formulas, including the orbit matrix used in the proof.

\paragraph{Structural conventions.}
Weakening, exchange, contraction, renaming, and ordinary-assumption cut are
admissible proof transformations. They are not additional checker
instructions. Substitution into a restricted derivation is used before the
bootstrap only when all substituted transport sources remain syntactically
nonzero. The main proof does not presuppose unrestricted substitution.

\subsection{Ordinary reading and immediate soundness}

On public formulas, read equality as actual numerical equality and
\[
 Q_yB\;:=\;B(0)\to\forall yB(y),\qquad
 D_{x,y}B\;:=\;\exists x\bigl(B(x,0)\to\forall yB(x,y)\bigr).
                                                               \tag{R}
\]
These readings are also sound for all the administrative rules in the
inherited scalar theory. Atomic DI follows by equation induction, compound
DI follows by numerical coverage, and DE is existential elimination. The
source-selected rule is an allowed equality or $Q$ elimination. Thus every
administrative proof maps to an old scalar proof under the ordinary reading.
What is \emph{not} asserted is completeness for that reading on every
administrative formula. The completeness and bi-interpretation claims have
the explicitly stated public domain $\CG$.

For one noncompleteness test later, a second sound reading of the
administrative language is useful: atomic $D$ means witnessed progress,
$\exists x\forall y(E(x,y)\to E(x,Sy))$; compound $D$ and every $Q$
keep the anchored reading in (R). This mixed reading agrees with (R) on
public guarded formulas. It is an auxiliary semantic argument, not an
extra primitive rule or an implicit change in the public interpretation.

% Research proof fragment. Written proof, not machine verification.
% Root may adapt into the final report. All claims are scoped below.
\section{The orbit bootstrap for guarded source dispatch}

The two existing reductions do not compose within the strictly guarded
calculus. They do compose when the proof language retains the administrative
orbit formulas described below. The distinction between public statements
and intermediate proof formulas is essential.

Let $\mathcal G$ be the existing guarded calculus. Its terms are
$x,0,St,I(n,a;u.h)$, and its public formulas are equations and
$D_{x,y}B$. An atomic matrix immediately below $D$ has the shape
\[
  I(y,0;u.h)=I(y,0;u.k),
\]
while a compound matrix is unrestricted within this recursive grammar.
Write $C(v,a,b)=I(v,a;u.b)$, with unused $u$, and
$G_v(s=t)\equiv C(v,0,s)=C(v,0,t)$.

For the intermediate argument, enlarge the grammar by also admitting the
\emph{orbit matrices}
\[
       O_a(y)\equiv y=I(y,a;u.Su),             \tag{O}
\]
where $a$ is independent of $y$ (it may contain the existential coordinate
and passive parameters). Call this enlarged grammar $\mathcal G^+$.
The primitive rules are those of the guarded calculus: computation, clause
use, $D$ elimination, and root-determined introduction (progress for atomic
matrices, cases for compound matrices). Restrict transport to a source term
that is not literally $0$. Call this restricted calculus $\mathcal R$.
No additional inference instruction is introduced by (O).

Contexts consist of reusable clauses $\langle A;y.B\rangle$, interpreted
as $A\to\forall yB$. An ordinary retained assumption is a constant clause,
as in the inherited checker. In what follows, all introduced $D$ matrices
are guarded, compound, or of form (O).

\subsection{The restriction that prevents a naive composition}

Call an equation \emph{zero-left} if its left term is literally $0$;
call all other equations and all compound formulas regular. If the source
$s$ of a transport is not literally zero, a zero-left minor cannot become
regular: if $p(s)$ is literally $0$, the numerical template $p$ must be the
constant $0$ template, so $p(t)$ is also literally $0$.

In the strictly guarded calculus, both endpoints of every atomic clause
opened by $D$ elimination are regular. Every atomic introduction premise
and its local assumption are also regular, since their left terms have an
$I$ head. Compound cases have compound endpoints. Consequently a retained
assumption $0=x$, for a fresh variable $x$, can be erased from every proof
of a regular conclusion in that calculus. A simultaneous induction on the
proof tree and its regular clause uses proves this assertion: a regular
transport conclusion has regular premises; a regular clause conclusion uses
a regular antecedent; the computation and binder rules preserve the
invariant. The only added zero-left constant clause cannot be used on a
path to a regular conclusion.

Thus $0=x\vdash x=0$ is not derivable there: otherwise $x=0$ would be a
theorem, contradicting the standard natural-number interpretation. This is
why one cannot simply quote the two previous admissibility theorems in
succession. In the present proof, orbit elimination has the nonregular
antecedent $0=I(0,a;u.Su)$ and can produce regular successor instances;
it is exactly the missing bridge.

\subsection{Safe equality and padding before zero-source symmetry}

Weakening, renaming, and ordinary-assumption cut are admissible. A numerical
substitution is admissible for a particular derivation if it leaves every
primitive transport source syntactically nonzero. No unrestricted
substitution into a restricted derivation is assumed.

For every $r$, let $b_r=I(0,r;u.0)$. The computation $b_r=r$, used as both
major and minor of transport with motive $v=r$, proves $r=r$.
From $p=q$ with $p\not\equiv0$, the same rule gives symmetry $q=p$ and
congruence $H(p)=H(q)$. Transitivity through a syntactically nonzero middle
term is immediate. Computation equations may therefore be reversed as
equations; using the reverse as a transport major still requires its new
source to be syntactically nonzero.

We also retain the inherited safe zero-join lemma:
\[
 p,q\not\equiv0,\quad p=0,\quad q=0 \quad\Longrightarrow\quad p=q.
                                                               \tag{ZJ}
\]
For completeness, put $z_0=I(0,0;u.0)$. From $b_p=b_p$ and $p=0$,
transport in $I(0,v;u.0)=b_p$ gives $z_0=b_p$. Reverse at source $z_0$,
then use $b_p=p$ at source $b_p$ to obtain $p=z_0$. Obtain $q=z_0$ in
the same way, reverse it at source $q$, and compose through $z_0$.
All sources are $p,q,z_0,b_p,b_q$, which are nonzero syntax.

\paragraph{A restricted lifting lemma.}
Suppose $p\not\equiv0$, and $a,b\not\equiv0$ have proved equations
$a=p$ and $b=q$. From $p=q$ one can derive $a=b$ without zero-source
transport. First compose through $p$ to get $a=q$. If $q\not\equiv0$,
reverse $b=q$ and compose through $q$. If $q\equiv0$, use (ZJ).
Conversely $a=b$ transports forward along $a=p$ and $b=q$ to $p=q$.

In particular, for every equation $E\equiv p=q$ whose left term is not
literally zero,
\[
     E\dashv\vdash G_{Sz}E,\qquad \vdash G_0E.
\]
The forward implication uses the lifting lemma with the two selector
computation equations. The reverse implication only contracts iterator
sources and in fact is valid for arbitrary $E$. Define
\[
       L(E)=D_{r,v}G_vE,
\]
with unused fresh $r,v$. For such a regular equation $E$, progress
introduction, ignoring its local hypothesis, proves $E\vdash L(E)$.
Elimination, the theorem $G_0E$, and use at $1$ prove $L(E)\vdash E$.

\paragraph{Numerical cases sufficient for the bootstrap.}
Cases for compound matrices need no transport: introduce a dummy-witness
$D$, open its clause, and use the supplied zero case at the target.
Consequently cases is also available for any equation template $E(y)$
whose left side is headed by $I$ or $S$. All its instances are regular,
so pad the two branches by $L$, apply compound cases, and unpad.
This restricted cases theorem suffices for every helper below. In
particular, we do not invoke cases with an arbitrary zero-left equation
as a retained assumption.

\subsection{Unconditional tests}

Define fixed finite terms
\[
\begin{aligned}
 Z(v)&=C(v,1,0),& K(v)&=C(v,0,1),& H(v)&=C(v,1,2),\\
 N(n)&=I(n,0;u.Z(u)),& R(n)&=I(n,0;u.H(u)),& L_1(n)&=N(R(n)).
\end{aligned}
\]
(The term $L_1$ is unrelated to the padding formula $L$.) The inherited
singleton-test calculation works with the restricted cases just proved:
\[
 H(v)=S K(v),\quad R(1)=1,\quad R(SSv)=2,\quad
 N(1)=1,\quad N(2)=0,
\]
\[
             L_1(1)=1,\qquad L_1(SSv)=0,\qquad
             L_1(Sv)=Z(v).                                \tag{T}
\]
Here are the details relevant to the restriction. For $H(v)=SK(v)$,
the zero and successor branches reduce to the nonzero numerals $1,2$;
join their endpoints through those numerals and apply restricted cases,
whose left side is $I$-headed. The values of $N$ and $R(1)$ are finite
forward computations. From
$R(Sv)=H(R(v))=SK(R(v))$, rewrite the argument of
$R(SSv)=H(R(Sv))$ using the $I$-headed source $R(Sv)$ and compute $H$ at
a successor to get $2$. This gives the values of $L_1$. For the last
equation in (T), the zero branch joins through $1$ and the successor
branch uses (ZJ) on the two $I$-headed terms that are zero. Its left side
$L_1(Sv)$ is again $I$-headed, so restricted cases applies.

For the filter, use cases on the \emph{reversed} equation
\[
                     C(Z(v),v,0)=v.                        \tag{F}
\]
Its left side is always $I$-headed. The zero branch is the forward
calculation $C(Z(0),0,0)=C(1,0,0)=0$. The successor branch is
$C(Z(Sv),Sv,0)=C(0,Sv,0)=Sv$. All replacements use iterator
sources. After cases has proved (F) at the arbitrary target $t$, safe
symmetry at its $I$-headed source gives
\[
                        t=C(Z(t),t,0).                    \tag{F'}
\]
All test lemmas are established in the empty context and only then
weakened into a context containing a zero-left assumption.

\subsection{The administrative orbit supplies the missing induction}

For a fixed term $t$ put $A_t(n)=I(n,t;u.Su)$. The formula
\[
                D_{r,y}\bigl(y=A_t(y)\bigr)               \tag{OC}
\]
with unused $r$ is derivable in $\mathcal R$, without a base hypothesis.
For progress introduction retain $z=A_t(z)$. Its source $z$ is a fresh
variable, hence syntactically nonzero. Successor congruence gives
$Sz=SA_t(z)$. Computation gives $A_t(Sz)=SA_t(z)$; reverse this equation
at its iterator source and compose through the successor term
$SA_t(z)$. Thus $Sz=A_t(Sz)$, exactly the required progress premise.
This derivation never substitutes $0$ for the variable-source transport.

Now retain $0=t$. If $t\equiv0$, symmetry is reflexivity. Otherwise
$t$ is a permitted source. Reverse the computation $A_t(0)=t$ and use
$t=A_t(0)$ at source $t$ in the assumption $0=t$ to get $0=A_t(0)$.
Open (OC), supply that base, and specialize its clause at $1$. This
gives $1=A_t(1)$. Forward computation gives $A_t(1)=St$, so
\[
                              1=St.
\]
Transport inside $L_1$ at source $1$, then contract the $I$-headed
sources using (T), obtaining $1=Z(t)$. Safe symmetry at source $1$
gives $Z(t)=1$. Apply this equation at its $I$-headed source to (F'),
then compute $C(1,t,0)=0$. The result is
\[
                       0=t\ \vdash_{\mathcal R}\ t=0.     \tag{ZS}
\]
No unrestricted numerical substitution or contextual cases principle
was used. The only new intermediate binder was (OC).

\subsection{Recovering the old guarded calculus}

From (ZS), safe symmetry, and (ZJ), obtain full symmetry, transitivity,
and numerical congruence exactly as in the existing source-dispatch
argument. For transitivity through literal zero, use (ZJ) if both
endpoints are nonzero syntax; otherwise one premise is already the
conclusion. For congruence along $0=t$, use reflexivity when $t$ is
literally $0$; otherwise reverse to $t=0$, apply safe congruence there,
and reverse the result. This derives unrestricted
transport in a literal equation by joining the two endpoint congruences.

Full transport through a $\mathcal G$ formula is now proved by external
induction on its binder count. In a compound-matrix $D$, use its case
introduction, transporting the base backwards, specializing the old
clause, and transporting the successor forwards by the smaller-formula
induction hypothesis. In a guarded atomic $D$, its source zero instance
is an unconditional theorem; hence the opened source clause supplies
its successor instance. Transport that literal equation forward, and
introduce the target by atomic progress, ignoring its local hypothesis.
For an orbit $D$ in $\mathcal G^+$, the target is already a theorem by
the derivation of (OC). This also establishes full transport throughout
$\mathcal G^+$ if that enlarged public grammar is desired.

Consequently every proof in the inherited guarded calculus translates
to $\mathcal R$, with the same public context and conclusion. All
nontransport rules are copied, and every old transport is replaced by
the preceding derived rule. The known translations and both formula
composites between the original scalar core and $\mathcal G$ therefore
remain available. The numerical domain, terms, and actual equality
are unchanged.

\subsection{Administrative soundness and the public theorem}

The compact combined checker admits more intermediate syntax than
$\mathcal G^+$. Use the already established mixed administrative reading:
an atomic $D$ matrix has the progressive meaning
$\exists x\forall y(E(x,y)\to E(x,Sy))$, and a compound matrix has the
anchored meaning $\exists x(B(x,0)\to\forall yB(x,y))$.
Clauses keep the anchored reading. Equation induction in iPRA makes
atomic elimination sound; introduction and the other rules are sound
directly. Restricting transport sources preserves this soundness.

On a public guarded atomic matrix, the base is provable, so the
progressive and anchored readings agree. On an orbit matrix, both
readings are unconditional theorems by the displayed progress proof
and equation induction; their agreement does not rely on a true base.
Thus both possible choices of public grammar, $\mathcal G$ and
$\mathcal G^+$, have the ordinary guarded meaning on their public
statements. Administrative proofs are conservative there by soundness
in iPRA and the completeness translation just proved.

The resulting precise statement is: for canonical guarded end sequents,
the combined administrative checker recognizes exactly the guarded
calculus's consequences, and its ordinary realization is equiconsistent
and whole-domain bi-interpretable with iPRA. This is not an equivalence
claim for every unrestricted administrative formula. It is also not
a claim that the compact checker validates the canonical public grammar;
if that is required, its parser must be included in the size accounting.
\section{Full interpretations and both composites}\label{sec:equivalence}

This section supplies the remaining link from the preceding orbit bootstrap
to the actual current native scalar core. It also records which portion of
the arithmetic bridge is inherited. Write $A\equider B$ for uniform
interderivability under arbitrary additional retained assumptions.

\subsection{Padding and arbitrary numerical cases after the bootstrap}

For an equation $E=(s=t)$ put
\[
 G_vE=(C(v,0,s)=C(v,0,t)),\qquad L(E)=D_{r,v}G_vE,
\]
where $r,v$ are fresh and $r$ is unused. Full equational transport, now
derived by the orbit bootstrap, proves
\[
 \vdash G_0E,\quad G_{Sz}E\equider E,\quad L(E)\equider E
                                                               \tag{P}
\]
for \emph{every} equation, including zero-left equations. Let
$P(A)=L(A)$ for an equation and $P(A)=A$ for a public compound formula.
Thus $P(A)$ is always compound and $P(A)\equider A$.

Compound numerical cases is already available from DI and DE. Applying it
to $L(E(y))$ gives numerical cases for arbitrary equations as well.
There is still a freshness condition: the context of the two case branches
must be independent of the case variable. In particular, one must not
perform cases on $y$ while retaining an unprocessed assumption $E(y)$.

\subsection{Contextual equation induction without an eigenvariable error}

We first specify the proof subsystem in which the following structural
operations are made. Let $\mathcal G^+$ be the guarded language enlarged
by the orbit matrices $y=I(y,a;u.Su)$, with $a$ independent of $y$.
Use these as object formulas, with their anchored $Q$ clauses as
administrative context resources. All proofs constructed in the orbit
bootstrap lie in this subsystem. Full transport has now been derived
throughout $\mathcal G^+$: the guarded and compound cases were proved
there, and an orbit $D$ target is already a theorem by its progress proof.

Temporarily allow this derived full-transport rule as an inference node.
Numerical substitution is admissible in that derived calculus by induction
on derivations, renaming eigenvariables first. The syntax of
$\mathcal G^+$ is preserved by substitution, the computation and binder
rules retain their forms, and full transport has no literal-source
restriction to invalidate. After substitution, replace every full-transport
node by its restricted derivation. This produces a derivation in $\UR$.
Thus substitutions into the constructed $\mathcal G^+$ proofs, including
their administrative orbit steps, are justified. This is a statement about
the specified clause subsystem, not about arbitrary administrative
derivations in the entire language of $\UR$.

For equations $U=(s=t)$ and $V=(r=q)$ define
\[
 M_w(U,V)=(C(w,s,r)=C(w,t,q)),\qquad
 U\Rightarrow_g V=D_{a,w}L(M_w(U,V)),
\]
with fresh unused $a$. The outer matrix is compound. The selector
computations and (P) identify its zero endpoint with $U$ and every
successor endpoint with $V$. Consequently this macro admits implication
introduction and elimination between equations. Introduction is compound
case DI; elimination opens its clause and uses it at $1$.

\begin{lemma}[Contextual equation induction]\label{lem:eqind}
For premise derivations in the $\mathcal G^+$ subsystem just described,
the following rule is admissible:
\[
 \frac{\Gamma\vdash E(0)\qquad\Gamma,E(z)\vdash E(Sz)}
      {\Gamma\vdash E(t)},
\]
where $E$ is a literal equation and $z$ is fresh for the fixed data.
\end{lemma}
\begin{proof}
Let $J(y)=G_yE(y)$. We seek the guarded progress premise
$\Gamma,J(z)\vdash J(Sz)$. Use cases on the \emph{compound} formula
\[
       \Phi(y)=J(y)\Rightarrow_g E(Sy)
\]
in the fixed context $\Gamma$. At zero, ignore $J(0)$, use the supplied
base $E(0)$ and the step at $0$. At a successor $Sw$, unguard $J(Sw)$ to
$E(Sw)$ and use the step at $Sw$. The two branches have no retained
assumption containing their case variable. Cases therefore gives
$\Gamma\vdash\Phi(z)$. Apply it under $J(z)$ and guard the result to
obtain $J(Sz)$.

Atomic DI yields $D_{r,y}J(y)$ with unused witness coordinate. Its base is
a theorem, so DE and specialization supply $J(Sw)$, hence $E(Sw)$.
Finally numerical cases for $E$, using the original base, yields $E(t)$.
The occurrences of the supplied step at $0$ and $Sw$ use substitution in
the derived full-transport $\mathcal G^+$ calculus, followed by expansion
of its transport nodes as specified above. No direct unsafe substitution
into a restricted primitive derivation is asserted.
\end{proof}

This lemma suffices for the recursive translation below: that translation
constructs only guarded object formulas, orbit intermediates, and their
anchored clauses. Its output therefore stays in $\mathcal G^+$, so every
recursive induction-step premise meets the stated hypothesis. After the
two interpretation theorems are proved, equation induction is also an
admissible rule of the full public consequence relation. It is not being
assumed in that broader form to establish those theorems.

\subsection{Translation of the current Q/D language}

Let $\NC$ be the current native scalar core. Its terms are the same terms,
and its formulas are equations, $Q_yB$, and $D_{x,y}B$.
Define $T$ recursively by
\[
 \begin{aligned}
 T(s=t)&=(s=t),\\
 T(Q_yB)&=D_{r,y}P(TB),\quad r\text{ fresh and unused},\\
 T(D_{x,y}B)&=D_{x,y}P(TB).
 \end{aligned}                                                   \tag{T}
\]
Capture-free renaming is understood. Every output belongs to $\CG$:
every translated matrix is compound, and the atomic matrices introduced
inside $L$ have the required zero guard. $T$ leaves every numerical term
unchanged and commutes with numerical substitution up to bound renaming.
Let $J:\CG\to\NC$ be literal inclusion.

\begin{proposition}[Both derivation translations]
For native $\Gamma,A$ and public $\Delta,B$,
\[
 \Gamma\vdash_{\NC}A\ \Longrightarrow\ T\Gamma\vdash_{\UR}TA,
 \qquad
 \Delta\vdash_{\UR}B\ \Longrightarrow\ J\Delta\vdash_{\NC}JB.
\]
\end{proposition}
\begin{proof}
For the second implication, read each administrative rule as the inherited
native rule. Its two introduction cases are allowed witnessed DI rules,
which the current native core derives; this is the upstream
\file{NativeInherited.lean} equivalence. Restricted transport is a native
derived rule. The administrative $Q$ is exactly native $Q$ for this direction.

For the first implication, translate by induction on the native derivation.
Computation and hypothesis are unchanged. Equational transport is available
from the orbit bootstrap. A translated $Q$ is a dummy-witness $D$ over
$P(TB)$; opening it supplies a clause whose endpoints are interderivable with
those of $TB$. This simulates QE by padding its base, using the clause at the
same term, and unpadding the result.

For DE, open the translated $D$ at a fresh witness $b$. Its clause has
antecedent $P(TB)(b,0)$ and consequent $P(TB)(b,y)$. Derive the entire
translated hypothesis
$T(Q_yB(b,y))=D_{r,y}P(TB(b,y))$ from this clause: introduce the
dummy-witness binder by compound case DI and, under its local base, use
the opened clause at $Sz$. Cut that formula into the recursively
translated minor, then close the outer elimination. This treats arbitrary
uses of the native $Q$ hypothesis, not merely its uses as a QE major.
The witness is absent from the unchanged final conclusion.
For native case DI, introduce the translated binder by compound case DI:
unpad its base, apply the translated premise, and pad its successor target.
For native direct equation progress, use Lemma~\ref{lem:eqind} under the
retained base to obtain every successor instance, then introduce the
dummy-witness translated $Q$ by compound case DI. These are effective
finite derivation transformations within the $\mathcal G^+$ clause
subsystem. Each use of structural substitution is made in its derived
full-transport calculus and then expanded, as specified above.
\end{proof}

\begin{proposition}[Both formula composites]\label{prop:composites}
For every native formula $A$ and public formula $B$,
\[
 A\equider_{\NC}JT(A),\qquad B\equider_{\UR}TJ(B).
\]
The equivalences are uniform under retained native or public contexts,
respectively.
\end{proposition}
\begin{proof}
In the native calculus, padding is equivalent to its argument. A
dummy-witness $D$ is equivalent to $Q$, and arbitrary native $D$ matrices
respect uniform interderivability. Structural induction proves the first
equivalence.

For the second, equations are unchanged. A compound public matrix respects
interderivability by opening the source binder, introducing the target by
case DI, converting its base backwards, using the source clause at a
successor, and converting forwards.
For a guarded atomic matrix $E(x,y)$, prove
$D_{x,y}E\equider D_{x,y}L(E)$. Forwards, open the source and use compound
case DI after unpadding the target base. Backwards, $E(b,0)$ is an
unconditional theorem; pad it, use the opened clause to obtain
$L(E(b,Sz))$, and unpad. Atomic progress DI introduces $D E$ while ignoring
its local induction hypothesis. Close DE. This establishes the atomic
case of the structural induction, and proves the second equivalence.
\end{proof}

\subsection{Composition with the inherited arithmetic bridge}

The current repository already proves whole-domain bi-interpretability of
the native scalar core and the usual primitive-recursive-function language
in Lean. The exact exported results are
\file{native_iff_pra}, \file{pra_iff_native},
\file{native_roundTrip}, \file{pra_native_roundTrip}, and
\file{native_inconsistent_iff} in
\file{Lean/ConciseIPRA/NativeBiInterpretation.lean} \cite{lean-native}.
The ordinary scalar theorem is \file{biInterp : BiInterp TeqAx PRAAx}
in \file{BiInterpretation.lean} \cite{lean-scalar}.

That bridge derives a predecessor, equality reflection into a numerical
term, quantifier-free induction, pairing, and compilation of primitive
recursion from scalar iteration. The exact attributed written arithmetic
sources and the two theorem-export files are included under
\file{reference/}. They are inherited dependencies, not discoveries in this
report. The existing Lean build was not rerun here, and these copied export
files are evidence of the upstream result, not a self-contained local Lean
project. The new orbit argument has not been formalized in Lean or Rocq.

Compose that bridge with $T$ and $J$. All term translations between the
public and native languages are literally identity maps. The inherited
scalar/iPRA term translations have internally proved identity composites.
Proposition~\ref{prop:composites} supplies the new formula composites.
Thus both resulting interpretations use the entire numerical domain and
actual equality, and both composite isomorphisms have graph $x=y$.
This proves Theorem~\ref{thm:main}, not merely mutual interpretability or
coincidence of true closed numerical equations.

The contradiction stays $0=1$. The public/native translations preserve it
literally, and the existing bridge identifies native inconsistency with
iPRA proving $\bot$. Equiconsistency follows as an external syntactic
theorem; no internal proof of iPRA's consistency is claimed.

\section{Complete checker sizes and the input contract}

\subsection{What is included in each measurement}

Each reported kernel is a complete standalone Python module importing only
the standard JSON parser. It reads one JSON certificate from standard input
and accepts exactly on normal exit with output \texttt{ACCEPT} followed by
a newline. Exceptions are rejection. No proof producer, translator, axiom
database, or helper is imported by a checker.

The universal recognition statement uses the idealized unbounded
integer-and-array semantics of the inherited sorted codec, together with
the present derivation translations. Ordinary Python has finite memory and
recursion limits, so it can reject sufficiently large valid certificates.
Successful executions retain the stated soundness guarantee. The finite
tests do not independently establish the inherited codec theorem.

We use the repository's normalization convention: remove genuine module,
function, and class docstrings, apply \texttt{ast.unparse}, append one
newline, and count significant tokenizer tokens, all AST nodes, and UTF-8
bytes. The runtime is CPython 3.12.14. Gzip level 9 with timestamp zero is
reported for information, not optimized as the primary objective.
The reference files were checked against their upstream Git blob hashes.

\begin{center}\small
\begin{tabular}{@{}lrrrr@{}}
\toprule
Complete program & Tokens & AST nodes & Bytes & Gzip-9\\
\midrule
Old AST profile & 263 & 283 & 607 & 348\\
New AST profile & 262 & \textbf{281} & 606 & 348\\
Old byte profile & 266 & 305 & 526 & 330\\
New byte profile & 265 & 303 & \textbf{525} & 331\\
Old token profile & 259 & 289 & 572 & 339\\
New token profile & \textbf{258} & 287 & 571 & 339\\
\bottomrule
\end{tabular}
\end{center}

The attained coordinate envelope therefore changes from
$(259,283,526)$ to $(258,281,525)$, in token/node/byte order.
Each triple combines coordinates from \emph{different programs}; no listed
program achieves all three. The byte-profile gzip grows by one byte.
Raw-source and normalized-source gzip measurements are both recorded in
the machine-readable metrics.

The dispatch change replaces the old independently selected mode test
\texttt{len(a) \% 2 == u} (or its token-profile counterpart) by
\texttt{not len(a) \% 2}. The source-dispatch restriction is retained.
The small source edit is justified by the substantial derivability argument
above; without the administrative orbit route, its apparent combination
would have the obstruction proved earlier.

\subsection{Canonical public inputs versus a strict interface}

The compact programs do not enforce the public guarded grammar.
Their theorem-recognition guarantee applies to its canonical encodings.
This convention must be part of the comparison, rather than an uncounted
front end added afterwards. Broader intermediate formulas are intentional,
and are indispensable to the new bootstrap.

\file{code/strict_checker.py} includes the same goal validator as the earlier
guarded package, coupled to the new kernel. It is a complete standalone
module and uses the unchanged public grammar.

\begin{center}\small
\begin{tabular}{@{}lrrrr@{}}
\toprule
Complete strict interface & Tokens & AST nodes & Bytes & Gzip-9\\
\midrule
Old scalar interface & 450 & 521 & 1109 & 538\\
Old guarded interface & 491 & 568 & 1208 & 574\\
New guarded interface & 491 & 563 & 1205 & 572\\
\bottomrule
\end{tabular}
\end{center}

The strict guarded interface improves by three bytes and five AST nodes
against its matched predecessor. It remains 96 bytes larger than the old
scalar interface. This is the principal practical limitation of calling
the overall presentation ``simpler''. The result improves the established
canonical-input frontier, not every possible complete-input contract.

\subsection{The entire 525-byte normalized module}

\begin{lstlisting}[language=Python]
import json

def c(a, g=[[]], k=0):
    i = lambda a, t=c, d=0: [i(b, t, d + 5 - len(a)) for b in a] if a * 0 == [] else a if k <= a < d else d % 2 and i(t, d=3)
    if i(a) in g:
        return
    m, s, u, t = q()
    a = m == ([s, u] if s else [t] * 3) and i(t, u) == a and i(t, s) or a
    h = c(m, g, k) or m * 3 or (u and [[s and [s], u, t], i(t, [s, u, t], 3) if s else u]) or (a := i(*a, t, 2))
    k -= 1
    c(i(a, [k]), g + [i(h, (m or not len(a) % 2) and k)], k)
q = json.load(open(0)).pop
c(q())
print('ACCEPT')
\end{lstlisting}

Here formula arities one, two, and three represent $D$, equality, and
administrative $Q$, respectively. Term arities one and three represent
successor and iteration. The opening operation uses the inherited sorted
arity codec and its dummy slots; the exact protocol is supplied in
\file{code/FORMAT.txt}, with the original codec justification under
\file{reference/}. The local-scope AST and byte programs and the
global-scope token program need separately numbered certificates.

The mode-free introduction does not remove the requirement that a witness
be chosen before its eigenvariable becomes available. The local proof
context must also be restored after checking a major premise. The new
programs preserve these scope conventions; they do not exchange a shorter
source for an unchecked freshness assumption.

\section{Validation, reproducibility, and evidence limits}

The archive contains a readable independent expansion, compact and strict
checkers, named proof producers, serialized certificates, rejection controls,
measurements, and source provenance. The executable evidence exercises the
orbit bootstrap and the recovered equation rules. There were 22 test cases
across five complete implementations, each run in ordinary Python and
\texttt{python -O}: all 220 expected outcomes matched. Fourteen cases are
positive canonical statements, one is a positive administrative orbit
statement (deliberately rejected by the strict public-goal interface), and
seven are negative controls. Four additional matched predecessor controls
also passed. The precise counts and verdicts are in
\file{evidence/validation.json} and
\file{evidence/matched_baseline_controls.json}. The producer is not trusted
by any kernel. Every generated proof is audited for primitive nonzero-source
transport and root-determined introduction before serialization.

The generic named derivation of $0=t\vdash t=0$ contains 326 distinct proof
nodes: 171 transport nodes, 122 computation nodes, and 10 introduction
nodes among them. Exactly one introduction is the unguarded orbit step,
and its progress hypothesis is used. The closed universal version has a
550-instruction local certificate. The independent guarded equation-induction
example uses its induction hypothesis and has 256 instructions.
Negative controls include the designated contradiction, an unallocated
closed parameter, an unavailable introduction witness, malformed records,
and an old atomic-case proof whose matrix is outside the public grammar.

The central distinction is between finite execution and universal proof.
The report gives the universal derivations and interpretation composites.
Certificates check concrete instances, scope behavior, and rejection of
invalid data. They are not a formal proof of every Python execution or a
Lean/Rocq formalization of the metatheorem.

Reproduction from the extracted directory requires Python's standard
library. Rebuilding the included PDF additionally requires a normal LaTeX
installation with the packages named in the source:
\begin{lstlisting}
python3 reproduce.py
pdflatex -interaction=nonstopmode -halt-on-error ipra_guarded_orbit.tex
pdflatex -interaction=nonstopmode -halt-on-error ipra_guarded_orbit.tex
\end{lstlisting}
No network fetch is required to run the supplied checker evidence or rebuild
the PDF. The report's LaTeX source is a single file. Frozen earlier theorem
exports are not claimed to build by themselves as a Lean project.

\subsection{A correction to the current README's substitution wording}

The current README says that substitution is ``not admissible'' in the
literal-zero-source restricted calculus. The precise issue is that a direct
rule-by-rule substitution can turn a permitted variable source into literal
zero. Admissibility itself follows from its already proved equivalence:
\[
 ZPrf\ \longrightarrow\ NPrf\ \xrightarrow{\ \mathrm{subst}\ }\
 NPrf\ \longrightarrow\ ZPrf.
\]
The first and last arrows are the two directions of
\file{zprf_iff_nprf}. This observation does not justify a circular bootstrap:
that equivalence must have been proved first. The new proof follows the same
discipline and makes no unrestricted substitution assumption before the
needed equality rules have been recovered.

\section{Further research questions}

\begin{enumerate}
\item \textbf{A smaller strict interface.} Can a single public syntax make
      the root-selected introduction complete while being simpler to
      validate than the current guarded grammar? The concrete next target
      is to beat the 1109-byte scalar strict comparator under the same
      complete-input contract.
\item \textbf{Mechanization.} Formalize the public/admin distinction and
      the stable-left cases lemma in Lean, then reuse the current
      \file{NativeZeroSource} arithmetic after replacing its general cases
      calls and its filter orientation. The orbit introduction should be
      the only genuinely new arithmetic bridge.
\item \textbf{Minimal intermediate language.} The completeness proof uses
      guarded formulas, their $Q$ clauses, and the orbit family
      $y=I(y,a;u.Su)$. Is there a comparably short checker that enforces this
      smaller intermediate grammar, with a simpler universal soundness
      proof than the full arity codec?
\item \textbf{Proof-size cost.} Quantify the expansion caused by safe
      zero-join, guarded padding, and the implication needed for contextual
      induction. A smaller kernel can have longer certificates; those are
      independent optimization questions.
\item \textbf{Uniform full-language introduction.} Can the independent
      case/progress choice be eliminated while retaining every original
      administrative formula and its ordinary reading, with no new
      validator cost? The present result establishes completeness only on
      the stated public language, so this stronger problem remains open.
\item \textbf{A mathematical gain beyond one-byte refinements.} Can a
      different numerical basis remove a constructor or a binding operation
      without replacing it by a richer arithmetic axiom or more complex
      parser? The prior register, tied-seed, and zero-start results should
      be treated as constraints, not rediscovered as new proposals.
\end{enumerate}

\section*{Attribution and claim ledger}
\addcontentsline{toc}{section}{Attribution and claim ledger}

The scalar term language, ordinary/native/iPRA bridge, source dispatch,
safe zero-join, singleton test, guarded padding, and raw certificate codec
come from the cited ProveIt research and formalization. The contribution
here is the failure analysis for strictly guarded composition and its repair
by administrative orbit formulas, including the restricted cases bootstrap
and the resulting compatible checker family. The equality-safe filter is
oriented in the reverse direction before the bootstrap.

The proved result is a bi-interpretable presentation and a small improvement
under an explicit comparison convention. No literature-wide priority,
global source-size lower bound, globally simplest axiomatization, shortest
proofs, new proof-assistant verification, or internal consistency proof is
claimed. The repository was read at a pinned commit; no remote changes were
published.

\begin{thebibliography}{9}
\bibitem{repo} Vladimir Reshetnikov, \emph{ProveIt: Concise presentations of
intuitionistic primitive recursive arithmetic}, current README at commit
\texttt{1517d9d79ce33f631e0f82960c849427cbba23bf}.
\url{https://github.com/VladimirReshetnikov/ProveIt/blob/1517d9d79ce33f631e0f82960c849427cbba23bf/Logic/PrimitiveRecursiveArithmetic/ConciseIPRA/README.md}.
\bibitem{source} ProveIt, \emph{A simpler elimination rule for intuitionistic
primitive recursive arithmetic}, \file{Research/source_dispatch/}, same
commit. Written proof, standalone checkers, and source provenance.
\url{https://github.com/VladimirReshetnikov/ProveIt/tree/1517d9d79ce33f631e0f82960c849427cbba23bf/Logic/PrimitiveRecursiveArithmetic/ConciseIPRA/Research/source_dispatch}.
\bibitem{guarded} ProveIt, \emph{A guarded core for intuitionistic primitive
recursive arithmetic}, \file{Research/guarded_core/}, same commit.
\url{https://github.com/VladimirReshetnikov/ProveIt/tree/1517d9d79ce33f631e0f82960c849427cbba23bf/Logic/PrimitiveRecursiveArithmetic/ConciseIPRA/Research/guarded_core}.
\bibitem{lean-native} ProveIt, \file{NativeBiInterpretation.lean}, same commit.
\url{https://github.com/VladimirReshetnikov/ProveIt/blob/1517d9d79ce33f631e0f82960c849427cbba23bf/Logic/PrimitiveRecursiveArithmetic/ConciseIPRA/Lean/ConciseIPRA/NativeBiInterpretation.lean}.
\bibitem{lean-scalar} ProveIt, \file{BiInterpretation.lean}, same commit.
\url{https://github.com/VladimirReshetnikov/ProveIt/blob/1517d9d79ce33f631e0f82960c849427cbba23bf/Logic/PrimitiveRecursiveArithmetic/ConciseIPRA/Lean/ConciseIPRA/BiInterpretation.lean}.
\bibitem{zero-lean} ProveIt, \file{NativeZeroSource.lean}, same commit.
\url{https://github.com/VladimirReshetnikov/ProveIt/blob/1517d9d79ce33f631e0f82960c849427cbba23bf/Logic/PrimitiveRecursiveArithmetic/ConciseIPRA/Lean/ConciseIPRA/NativeZeroSource.lean}.
\bibitem{arithmetic} ProveIt, attributed scalar arithmetic proofs in
\file{Research/certified_zero/reference/}, same commit:
\file{inherited_choice_basis.tex}, \file{inherited_ordinary_bridge.tex},
\file{scalar_predecessor.tex}, \file{equality_reflection.tex}, and
\file{primitive_recursion.tex}. Copies are supplied in this archive under
the repository's MIT No Attribution license.
\end{thebibliography}
\end{document}
